% ---------------------------------------------------------------
% Plantilla de avance -- Matematicas Aplicadas al Modelado de Materiales
%
% Se entrega el .tex y el .pdf compilado, en:
%   alumnos/<codigo>/avances/tNN/
% junto con las carpetas codigo/, datos/ y figuras/
%
% Compilar con:  pdflatex avance.tex
% ---------------------------------------------------------------
\documentclass[11pt,letterpaper]{article}

\usepackage[utf8]{inputenc}
\usepackage[T1]{fontenc}
% La primera vez, MikTeX descarga solo el paquete de espanol. Si le pregunta,
% acepte la instalacion.
\usepackage[spanish,mexico]{babel}
\usepackage[margin=2.5cm]{geometry}
\usepackage{amsmath,amssymb}
\usepackage{graphicx}
\usepackage{booktabs}
\usepackage{hyperref}

\graphicspath{{figuras/}}

\title{Avance tNN: \textit{titulo corto de lo que hizo}}
\author{Nombre del estudiante \\ Codigo: 000000000}
\date{\today}

\begin{document}

\maketitle

% ===============================================================
\section{Encargo recibido}
% Copiado literal del cierre del turno anterior, sin reescribirlo.
% Las cuatro partes: concepto, fuente, pregunta que debe responder,
% evidencia esperada.
% ===============================================================

\begin{description}
  \item[Concepto o tecnica:]
  \item[Fuente:]
  \item[Pregunta que debo poder responder:]
  \item[Evidencia esperada:]
\end{description}

% ===============================================================
\section{Que estudie}
% El concepto con SUS palabras. Toda ecuacion lleva debajo la lista
% de sus variables con unidades. No copie el libro: explique.
% ===============================================================

\begin{equation}
  y = f(x)
\end{equation}

donde:

\begin{itemize}
  \item $y$: descripcion (unidades)
  \item $x$: descripcion (unidades)
\end{itemize}

% ===============================================================
\section{Como lo aplique a mi caso}
% La tecnica sobre SU fenomeno, con sus numeros. Si esta seccion
% se puede leer sin saber cual es su tesis, esta mal escrita.
% ===============================================================

% ===============================================================
\section{Evidencia}
% Figuras numeradas y referenciadas desde el texto, tabla de
% resultados y el error cuantificado con sus unidades.
% ===============================================================

\begin{figure}[htbp]
  \centering
  % \includegraphics[width=0.75\textwidth]{nombre_figura.png}
  \caption{Que se ve aqui y que hay que concluir de ello.}
  \label{fig:resultado}
\end{figure}

\begin{table}[htbp]
  \centering
  \caption{Metricas de error del modelo ajustado.}
  \label{tab:metricas}
  \begin{tabular}{lrl}
    \toprule
    Metrica & Valor & Unidades \\
    \midrule
    MAE  & & \\
    RMSE & & \\
    $R^2$ & & (adimensional) \\
    \bottomrule
  \end{tabular}
\end{table}

% ===============================================================
\section{Como reproducirlo}
% El panel va a seguir esto al pie de la letra. Que script se corre,
% en que orden, y que archivo produce cada uno.
% ===============================================================

\begin{enumerate}
  \item \texttt{codigo/gen\_datos.m} produce \texttt{datos/datos.csv}
  \item \texttt{codigo/ajuste.m} lee \texttt{datos/datos.csv} y produce
        \texttt{figuras/ajuste.png} y \texttt{datos/metricas.csv}
\end{enumerate}

Version de Matlab utilizada:

Toolboxes necesarias (o la alternativa sin toolbox):

% ===============================================================
\section{Que no me salio}
% Lo que fallo y lo que quedo abierto. Un avance sin esta seccion
% esta incompleto. Reportar el problema no baja la calificacion.
% ===============================================================

\begin{itemize}
  \item
\end{itemize}

\subsection*{Lo que voy a preguntar en la sesion}

% ===============================================================
\section{Referencias}
% ===============================================================

\begin{thebibliography}{9}

\bibitem{fuente}
  Autor, \textit{Titulo}, editorial, ano, capitulo y seccion.

\end{thebibliography}

\end{document}
